
\documentclass{article}
\usepackage{colortbl}
\usepackage{makecell}
\usepackage{multirow}
\usepackage{supertabular}

\begin{document}

\newcounter{utterance}

\centering \large Interaction Transcript for game `mmlu\_pro', experiment `mmlu\_pro\_physics', episode 406 with Qwen3.6{-}27B.
\vspace{24pt}

{ \footnotesize  \setcounter{utterance}{1}
\setlength{\tabcolsep}{0pt}
\begin{supertabular}{c@{$\;$}|p{.15\linewidth}@{}p{.15\linewidth}p{.15\linewidth}p{.15\linewidth}p{.15\linewidth}p{.15\linewidth}}
    \# & \multicolumn{2}{c}{Player} && \multicolumn{2}{c}{Game Master} \\
    \hline

    \theutterance \stepcounter{utterance}  
    & & & \multicolumn{4}{p{0.6\linewidth}}{
        \cellcolor[rgb]{0.9,0.9,0.9}{
            \makecell[{{p{\linewidth}}}]{
                \texttt{\tiny{[P1$\langle$GM]}}
                \texttt{The following are multiple choice questions (with answers) about physics. Answer only with "The answer is (X)" where X is the correct letter choice.} \\
\\ 
\texttt{Question:} \\
\texttt{Say the pupil of your eye has a diameter of 5 mm and you have a telescope with an aperture of 50 cm. How much more light can the telescope gather than your eye?} \\
\texttt{Options:} \\
\texttt{A. 1000 times more} \\
\texttt{B. 50 times more} \\
\texttt{C. 5000 times more} \\
\texttt{D. 500 times more} \\
\texttt{E. 10000 times more} \\
\texttt{F. 20000 times more} \\
\texttt{G. 2000 times more} \\
\texttt{H. 100 times more} \\
\texttt{I. 10 times more} \\
\texttt{J. N/A} \\
\texttt{The answer is (E)} \\
\\ 
\texttt{Question:} \\
\texttt{Where do most short{-}period comets come from and how do we know?} \\
\texttt{Options:} \\
\texttt{A. The Kuiper belt; short period comets tend to be in the plane of the solar system just like the Kuiper belt.} \\
\texttt{B. The asteroid belt; short period comets tend to come from random directions indicating a spherical distribution of comets called the asteroid belt.} \\
\texttt{C. The asteroid belt; short period comets tend to be in the plane of the solar system just like the asteroid belt.} \\
\texttt{D. The Oort cloud; short period comets have orbital periods similar to asteroids like Vesta and are found in the plane of the solar system just like the Oort cloud.} \\
\texttt{E. The Oort Cloud; short period comets tend to come from random directions indicating a spherical distribution of comets called the Oort Cloud.} \\
\texttt{F. The Oort cloud; short period comets tend to be in the plane of the solar system just like the Oort cloud.} \\
\texttt{G. The asteroid belt; short period comets have orbital periods similar to asteroids like Vesta and are found in the plane of the solar system just like the asteroid belt.} \\
\texttt{H. N/A} \\
\texttt{I. N/A} \\
\texttt{J. N/A} \\
\texttt{The answer is (A)} \\
\\ 
\texttt{Question:} \\
\texttt{A refracting telescope consists of two converging lenses separated by 100 cm. The eye{-}piece lens has a focal length of 20 cm. The angular magnification of the telescope is} \\
\texttt{Options:} \\
\texttt{A. 10} \\
\texttt{B. 40} \\
\texttt{C. 6} \\
\texttt{D. 25} \\
\texttt{E. 15} \\
\texttt{F. 50} \\
\texttt{G. 30} \\
\texttt{H. 4} \\
\texttt{I. 5} \\
\texttt{J. 20} \\
\texttt{The answer is (H)} \\
\\ 
\texttt{Question:} \\
\texttt{Colors in a soap bubble result from light} \\
\texttt{Options:} \\
\texttt{A. dispersion} \\
\texttt{B. deflection} \\
\texttt{C. refraction} \\
\texttt{D. reflection} \\
\texttt{E. interference} \\
\texttt{F. converted to a different frequency} \\
\texttt{G. polarization} \\
\texttt{H. absorption} \\
\texttt{I. diffraction} \\
\texttt{J. transmission} \\
\texttt{The answer is (E)} \\
\\ 
\texttt{Question:} \\
\texttt{A microwave oven is connected to an outlet, 120 V, and draws a current of 2 amps. At what rate is energy being used by the microwave oven?} \\
\texttt{Options:} \\
\texttt{A. 240 W} \\
\texttt{B. 120 W} \\
\texttt{C. 10 W} \\
\texttt{D. 480 W} \\
\texttt{E. 360 W} \\
\texttt{F. 200 W} \\
\texttt{G. 30 W} \\
\texttt{H. 150 W} \\
\texttt{I. 60 W} \\
\texttt{J. 300 W} \\
\texttt{The answer is (A)} \\
\\ 
\texttt{Question:} \\
\texttt{A particle of mass \$m\$ and velocity \$u\_1\$ makes a head{-}on collision with another particle of mass \$2 m\$ at rest. If the coefficient of restitution is such to make the loss of total kinetic energy a maximum, what are the velocities \$v\_1\$ after the collision?} \\
\texttt{Options:} \\
\texttt{A. \$0\$} \\
\texttt{B.   \$\textbackslash{}frac\{1\}\{3\}\$\$u\_1\$} \\
\texttt{C. \$\textbackslash{}frac\{1\}\{2\}\$\$u\_1\$} \\
\texttt{D. \$\textbackslash{}frac\{2\}\{5\}u\_1\$} \\
\texttt{E. \$\textbackslash{}frac\{2\}\{3\}\$\$u\_1\$} \\
\texttt{F. \$\textbackslash{}frac\{3\}\{4\}u\_1\$} \\
\texttt{G. \${-}\textbackslash{}frac\{1\}\{2\}u\_1\$} \\
\texttt{H. \${-}\textbackslash{}frac\{1\}\{3\}u\_1\$} \\
\texttt{I. \$u\_1\$} \\
\texttt{J. \$\textbackslash{}frac\{1\}\{4\}u\_1\$} \\
            }
        }
    }
     \\ \\

    \theutterance \stepcounter{utterance}  
    & \multicolumn{4}{p{0.6\linewidth}}{
        \cellcolor[rgb]{0.561,0.737,0.561}{
            \makecell[{{p{\linewidth}}}]{
                \texttt{\tiny{[P1$\rangle$GM]}}
                \texttt{The answer is (H)} \\
            }
        }
    }
    & & \\ \\

    \theutterance \stepcounter{utterance}  
    & & & \multicolumn{2}{p{0.3\linewidth}}{
        \cellcolor[rgb]{0.9,0.9,0.9}{
            \makecell[{{p{\linewidth}}}]{
                \texttt{\tiny{[GM$|$GM]}}
                \texttt{H} \\
            }
        }
    }
    & & \\ \\

    \theutterance \stepcounter{utterance}  
    & & & \multicolumn{2}{p{0.3\linewidth}}{
        \cellcolor[rgb]{0.9,0.9,0.9}{
            \makecell[{{p{\linewidth}}}]{
                \texttt{\tiny{[GM$|$GM]}}
                \texttt{B} \\
            }
        }
    }
    & & \\ \\

    \theutterance \stepcounter{utterance}  
    & & & \multicolumn{2}{p{0.3\linewidth}}{
        \cellcolor[rgb]{0.9,0.9,0.9}{
            \makecell[{{p{\linewidth}}}]{
                \texttt{\tiny{[GM$|$GM]}}
                \texttt{game\_result = LOSE} \\
            }
        }
    }
    & & \\ \\

\end{supertabular}
}

\end{document}
